\subsection{Defined Experiments}
\label{sec:budget}
\textbf{E1: diagnosis.} We run 27 cells in three cyclic blocks with the same
four-core budget. They include ingestion alone and two existing supported
prompts under inference alone, shared execution, a fixed 2+2 split and a fixed
3+1 split. Prefill and decode thread counts follow the LLM allocation. The
encoder receives two threads in 2+2 and one in 3+1. We set affinity before
creating threads and check it during each run. Warmup and v1 index preparation
are outside timing. Prompt KV reuse is off. These policies change both thread
counts and CPU placement, so they do not isolate cache or bandwidth effects.

\textbf{E2: application.} We run 36 episodes with three added articles, two
arrival offsets (0/90\,s), three policies (shared, 2+2, deferred) and two ways
of updating the index. Incremental publication embeds 82 new chunks. A rebuild
embeds the same complete 269-chunk v2 corpus in an empty staging index, then
switches the active pointer atomically. The old v1 remains readable during
the build. We do not pad the corpus to make updates slower.
The new request arrives at its scheduled time, waits for v2, retrieves the
top three passages and submits to the server. The old request may still occupy
the only inference slot. We record submission, first token and engine phase
timers. The engine does not expose the exact time a request leaves its queue.
The old request uses the same supplied context that passed the earlier checks.

\textbf{E3: timing test.} We use the next three articles, arrivals at 30/60\,s
and the two concurrent policies, giving 12 episodes. Before running them, we
use E2 incremental medians for each policy to predict current completion,
publication, retrieval and new-request service time. Engine prompt and
generation timers exclude queue waiting. The completion estimate is
$\max\{T_c,\max(T_a,T_p)+T_r\}+T_s$, where $T_a$ is arrival, $T_r$ is retrieval
and $T_s$ is service time. We save hashes of the parameters and predictions
and do not fit them again. These articles have appeared in earlier model
runs. They test timing on other cases, not semantic accuracy on unseen data.
An infrastructure failure interrupted E2 after 17 saved episodes. We retained
those episodes and ran the remaining 19 with the same experimental settings.
